\section*{الفصل \ref{ch:b1:chemical-biogenic-sediments} --- الرسوبيات الكربوناتية والسيليسية والتبخّرية}

\begin{solution}{exo:b1:chemical-biogenic-sediments:1}
الكلوريد $0.5503\times35.165 = \qty{19.35}{g}$؛ الصوديوم $0.3066\times35.165 =
\qty{10.78}{g}$؛ الكبريتات \qty{2.711}{g}؛ الماغنيسيوم \qty{1.284}{g}. ويكوّن الصوديوم والكلوريد معًا $0.5503 + 0.3066 = 0.857$، أي \qty{85.7}{\%} من الملح.
\end{solution}

\begin{solution}{exo:b1:chemical-biogenic-sediments:2}
الكسر المتبقّي هو $1/\mathrm{CF}$: كربونات الكالسيوم $1/1.8 = 0.556$ متبقّية، أي \qty{44.4}{\%} تبخّرت؛ الجبس $0.263$ متبقٍّ، \qty{73.7}{\%} تبخّرت؛ الهاليت $0.094$ متبقٍّ، \qty{90.6}{\%} تبخّرت؛ كبريتات الماغنيسيوم $0.014$ متبقٍّ، \qty{98.6}{\%} تبخّرت.
\end{solution}

\begin{solution}{exo:b1:chemical-biogenic-sediments:3}
(أ)~$\rho_s(1-\phi) = G/a = 8.4/0.0042 = \qty{2000}{kg.m^{-3}}$. (ب)~مع الأراغونيت، $1 - 2000/2930 = 0.32$، أي مسامية \qty{32}{\%}؛ ومع الكالسيت، $1 - 2000/2710 = 0.26$. (ج)~$a = 3.7/2000 = \qty{1.85e-3}{m}$، أي \qty{1.85}{mm} في السنة.
\end{solution}

\begin{solution}{exo:b1:chemical-biogenic-sediments:4}
حتى \qty{4000}{m} ينخفض المحتوى ببطء فقط (من 88 إلى 80)؛ وتحت \qty{4000}{m} ينخفض بحدّة: \omterm{def:b1:chemical-biogenic-sediments:ccd}{فالحدّ الانحلالي} يقع عند نحو \qty{4000}{m}. ويهبط المحتوى دون \qty{10}{\%} بين \qty{4700}{m} (12) و\qty{4800}{m} (6): $4700 + 100\times(12-10)/(12-6) = \qty{4733}{m}$، فالتعويض الكربوناتي عند نحو \qty{4730}{m}، أي نحو \qty{700}{m} تحت \omterm{def:b1:chemical-biogenic-sediments:ccd}{الحدّ الانحلالي}.
\end{solution}

\begin{solution}{exo:b1:chemical-biogenic-sediments:5}
(أ)~\ce{Ca^2+ + 2 HCO3- -> CaCO3 + CO2 + H2O}: مول واحد من \ce{CO2} لكل مول من \ce{CaCO3}، ف$44.01/100.09 = \qty{0.440}{kg}$ من \ce{CO2} لكل كيلوغرام من \ce{CaCO3}. (ب)~$3.7\times0.440 = \qty{1.63}{kg}$ من \ce{CO2} لكل متر مربّع في السنة. (ج)~من أيوني بيكربونات الهيدروجين المستعملين، لا يدخل إلى \omterm{def:b1:rock-forming-minerals:mineral}{المعدن} إلا كربون واحد؛ والآخر يخرج ثاني أكسيد الكربون. والتكلّس ينزع القلوية من الماء ويخفّض قدرته على حمل الكربون المذاب؛ وما يصل من \ce{CO2} إلى الهواء يتوقّف على تنظيم ماء البحر، وهذه مسألة توازنات الكربونات في كتب الكيمياء.
\end{solution}

\begin{solution}{exo:b1:chemical-biogenic-sediments:6}
(أ)~$t_c = (2000/350)^2 = \qty{32.7}{Myr}$؛ بمعدّل \qty{25}{km} لكل مليون سنة، $25\times32.7 = \qty{816}{km}$ من المحور. (ب)~$t_c = (2600/350)^2 =
\qty{55.2}{Myr}$، أي \qty{1380}{km}. (ج)~القاع بعمر \qty{60}{Myr} تلقّى الطين الأحمر لمدّة $60 - 32.7 = \qty{27.3}{Myr}$ في الحالة الأولى، ولمدّة $60 - 55.2 = \qty{4.8}{Myr}$ فقط في الثانية.
\end{solution}

\begin{solution}{exo:b1:chemical-biogenic-sediments:7}
(أ)~نعم: $a = 4.2 > r = 3$، والشعبة تكسب \qty{1.2}{mm} في السنة على البحر حتى تبلغ قمّتها السطح. (ب)~ينمو $D$ بمقدار $10 - 4.2 = \qty{5.8}{mm}$ في السنة؛ ومن \qty{2}{m} إلى \qty{50}{m} يستغرق ذلك $48/0.0058 = \num{8300}$ سنة. (ج)~لأن معدّل نموّها يهبط كلما زاد عمق الماء: فلا يصل إلى الطحالب التي تعيش في المرجانيات وتغذّيها إلا قليل من الضوء، فينخفض $a$، والشعبة تتأخّر أسرع فأسرع.
\end{solution}

\begin{solution}{exo:b1:chemical-biogenic-sediments:8}
(أ)~التفاعل هو \ce{CaSO4.2H2O -> CaSO4 + 2 H2O}. (ب)~الحجوم المولارية هي $172.17/2.30 =
\qty{74.9}{cm^3}$ و$136.14/2.97 = \qty{45.8}{cm^3}$، فالصلب يفقد $1 - 45.8/74.9 = \qty{38.8}{\%}$ من حجمه. (ج)~$2\times18.02/172.17 = 0.209$: \qty{20.9}{\%} من الكتلة يخرج ماءً. وحين يُطلق هذا الماء في العمق يرفع الضغط في مسامات الصخور المحيطة؛ وتحت غطاء محكم يمكن أن يُضعف الصخور ويساعد الصدوع على الانزلاق، ويمكنه أن يُذيب أملاحًا أخرى وينقلها.
\end{solution}

\begin{solution}{exo:b1:chemical-biogenic-sediments:9}
(أ)~$9.2/255 = \qty{3.6}{\%}$ محفوظ، و\qty{96.4}{\%} منحلّ. (ب)~حيث يكون الإنتاج عاليًا جدًّا، تبقى بضعة في المئة منه تدفّقًا كبيرًا، وحيث يخفّفه قليل من الطين أو الكربونات: تحت الماء الصاعد في المحيط الجنوبي، وفي شمال المحيط الهادئ، وفي المحيط الهادئ الاستوائي، وتحت التعويض الكربوناتي. (ج)~الهياكل الرقيقة الهشّة تذوب أولًا وتبقى الأقوى: فيحفظ الرسوب نخبة من أقوى الأنواع، لا المجتمع الذي عاش في الماء.
\end{solution}

\begin{solution}{exo:b1:chemical-biogenic-sediments:10}
(أ)~$EA = 1.5\times\qty{1e10}{m^2} = \qty{1.5e10}{m^3}$ في السنة. (ب)~مع $Q_{\mathrm{in}} = \qty{1.8e10}{m^3}$، و$Q_{\mathrm{out}} = \qty{0.3e10}{m^3}$، و$k = 6$: يتجاوز عتبة الجبس (3.8) دون عتبة الهاليت (10.6)، فيترسّب الكربونات والجبس. ومع $Q_{\mathrm{in}} =
\qty{1.575e10}{m^3}$، و$Q_{\mathrm{out}} = \qty{0.075e10}{m^3}$، و$k = 21$، لأيون يبقى في المحلول مثل الماغنيسيوم: يجتاز المحلول عتبة الهاليت، فيترسّب الهاليت أيضًا، لكن لا أملاح الماغنيسيوم (70). (ج)~يقتضي الهاليت أن يكون $Q_{\mathrm{in}}/(Q_{\mathrm{in}} - EA) \ge 10.6$، أي $Q_{\mathrm{in}} \le 10.6\,EA/9.6 = 1.104\,EA = \qty{1.66e10}{m^3}$ في السنة: فقد يتجاوز الداخل التبخّر بعشر تقريبًا على الأكثر.
\end{solution}

\begin{solution}{exo:b1:chemical-biogenic-sediments:11}
(أ)~يحتوي ماء البحر على \qty{67.3}{mg} من البروميد لكل كيلوغرام؛ وعند \num{10.6} مرة، يحتوي المحلول على $67.3\times10.6 = \qty{713}{mg.kg^{-1}}$، وأول هاليت $0.032\times713 = \qty{22.8}{mg.kg^{-1}}$. (ب)~عند 30 مرة، $67.3\times30 =
\qty{2019}{mg.kg^{-1}}$ في المحلول و\qty{64.6}{mg.kg^{-1}} في الهاليت. (ج)~لا يأخذ الهاليت إلا \num{3.2} من مئة من تركيز المحلول، فيبقى البروميد في المحلول تقريبًا كلّه ويتركّز أكثر مع تبخّر الماء. (د)~الصعود يسجّل تركيز المحلول؛ والهبوط المفاجئ يسجّل تخفيفه بتدفّق جديد من ماء البحر إلى الحوض، أو بماء عذب، قبل أن يواصل الهاليت نموّه.
\end{solution}

\begin{solution}{exo:b1:chemical-biogenic-sediments:12}
(أ)~$t_c = \qty{32.65}{Myr}$، ف$0.02\times\num{32650} = \qty{653}{m}$ من الرسوب الكلسي (\num{32650} فترة من ألف سنة بسمك \qty{2}{cm}) و$0.001\times\num{37350} = \qty{37}{m}$ من الطين الأحمر. (ب)~$t_c = (2500/350)^2 =
\qty{51.0}{Myr}$: \qty{1020}{m} من الرسوب الكلسي و\qty{19}{m} من الطين الأحمر. فالتعويض الأعمق يعطي رسوبًا كلسيًّا أسمك بكثير، والطين الأحمر في قاع الأطلسي العميق أرقّ منه في المحيط الهادئ للعمر نفسه. (ج)~ربما كان التعويض الكربوناتي أعمق لفترة، فصار الموقع فوقه مرة أخرى؛ أو حملت الصفيحة الموقع تحت حزام إنتاجية عالية، كخطّ الاستواء، حيث سقط كربونات أكثر؛ أو نُقلت الطبقة بانزلاق أو بتيار عكارة من جزء أعلى من قاع البحر.
\end{solution}

\begin{solution}{pb:b1:chemical-biogenic-sediments:1}
\textbf{1.} $1030\times1000 = \qty{1.03e6}{kg}$ من ماء البحر، يحمل
$1.03\times10^6\times0.035165 = \qty{36200}{kg}$ من الملح.

\textbf{2.} الجبس $1.03\times10^6\times1.581\times10^{-3} = \qty{1628}{kg}$؛ الهاليت $1.03\times10^6\times0.02741 = \qty{28230}{kg}$.

\textbf{3.} كربونات الكالسيوم $113.3/2710 = \qty{0.042}{m}$؛ الجبس $1628/2300 =
\qty{0.71}{m}$؛ الهاليت $28230/2170 = \qty{13.0}{m}$.

\textbf{4.} للأنهيدريت الكتلة $1628\times136.14/172.17 =
\qty{1288}{kg}$ والسماكة $1288/2970 = \qty{0.43}{m}$.

\textbf{5.} جافًّا، الملح $0.110 + 1.250 + 27.41 + 6.34 = \qty{35.11}{g}$ لكل كيلوغرام، منه الهاليت $27.41/35.11 = \qty{78}{\%}$: يكوّن الصوديوم والكلوريد معظم الملح المذاب، ويأخذ الهاليت كلّهما تقريبًا.

\textbf{6.} من الأسفل إلى الأعلى: كربونات الكالسيوم، ثم الجبس، ثم الهاليت، ثم أملاح المرّ. يبدأ كل \omterm{def:b1:rock-forming-minerals:mineral}{معدن} بالترسّب حين يبلغ المحلول تشبّعه، والأقلّ ذوبانًا أولًا، ويستقرّ فوق ما ترسّب قبله.

\textbf{7.} عمق المحلول هو $1000/\mathrm{CF}$: \qty{556}{m} لكربونات الكالسيوم، و\qty{263}{m} للجبس، و\qty{94}{m} للهاليت.

\textbf{8.} $1 - 1/10.6 = \qty{90.6}{\%}$.

\textbf{9.} $1000/70 = \qty{14}{m}$ لكبريتات الماغنيسيوم، و$1000/90 = \qty{11}{m}$ لأملاح البوتاسيوم.

\textbf{10.} نحو \qty{1.10}{g.cm^{-3}} حين تتكوّن كربونات الكالسيوم، و\num{1.13} مع الجبس، و\num{1.22} مع الهاليت، وأكثر من \num{1.30} مع أملاح المرّ، مقابل \num{1.03} لماء البحر. يطفو ماء البحر الداخل فوق المحلول الكثيف، الذي يبقى في القاع؛ وفي حوض محجوز يمكن للمحلول أن يعود إلى المحيط في العمق تحت الداخل.

\textbf{11.} لا يبلغ سمكها إلا \qty{4}{cm}، مقابل \qty{13}{m} من الهاليت؛ وتتكوّن على الهوامش، حيث يتركّز الماء الداخل أولًا، أكثر منه في الوسط، فتختلط بقاعدة الجبس أو يحلّ محلّها.

\textbf{12.} يأخذ الهاليت جزءًا صغيرًا فقط من بروميد المحلول (التمرين~\ref{exo:b1:chemical-biogenic-sediments:11})، فيصير المحلول أغنى في البروميد كلما نما الهاليت، وتأخذ كل طبقة جديدة من الهاليت منه أكثر.

\textbf{13.} $1500/13.01 = 115.3$: 116 ملءً.

\textbf{14.} $116\times0.434 = \qty{50}{m}$ من الأنهيدريت، لا كطبقة واحدة بل ك\num{116} طبقة، كلّ منها أقلّ قليلًا من نصف متر سمكًا، عند قاعدة الهاليت من ملئها: تعاقب دوري من الأنهيدريت والهاليت.

\textbf{15.} لكل ملء، $1.03\times10^6\times6.34\times10^{-3}/2000 =
\qty{3.3}{m}$، ف\qty{379}{m} ل\num{116} ملءً. وهي أكثر الأملاح ذوبانًا: فالداخل التالي من ماء البحر يعيد إذابتها، والحوض الذي يخرج محلوله قبل النهاية لا يصل إليها أبدًا.

\textbf{16.} عند كل إعادة ملء يُخفَّف المحلول: فينخفض محتوى البروميد بحدّة عند قاعدة الهاليت الجديد ثم يرتفع من جديد نحو الأعلى، في نمط مُسنَّن بسنّة واحدة لكل ملء.

\textbf{17.} $2000/13.01 = 154$ ملءً.

\textbf{18.} $5000\times365.25\times86400 = \qty{1.578e11}{m^3}$ في السنة، موزّعةً على $\qty{1.0e11}{m^2}$: $E = \qty{1.58}{m}$ في السنة.

\textbf{19.} $1.578\times10^{11}\times1030\times0.035165 = \qty{5.72e12}{kg}$ من الملح في السنة.

\textbf{20.} $1.578\times10^{11}\times1030\times0.02741/2170 =
\qty{2.05e9}{m^3}$ من الهاليت في السنة، أي \qty{0.0205}{m}، نحو \qty{2}{cm} في السنة على القاع كله.

\textbf{21.} $1500/0.0205 = \num{73000}$ سنة.

\textbf{22.} ينخفض المستوى بمقدار \qty{1.58}{m} في السنة؛ و\qty{1000}{m} من الماء تذهب في $1000/1.578 = 634$ سنة.

\textbf{23.} يظهر الهاليت حين يذهب \qty{90.6}{\%} من الماء، بعد $905.7/1.578 = 574$ سنة.

\textbf{24.} استغرق \qty{1.5}{km} من الهاليت ملح 116 ملءً للحوض (\num{115.3}، يُقرَّب إلى الأعلى)، سواء جاءت الواحدة بعد الأخرى مع جفاف الحوض وإعادة ملئه، أو بشكل متّصل فوق العتبة، كما في السؤال~21، بمعدّل حوض واحد من ماء البحر كل $1000/1.578 = 634$ سنة.
\end{solution}
